\documentclass[11pt,a4paper]{article}

\usepackage[margin=2.4cm]{geometry}
\usepackage{amsmath}
\usepackage{booktabs}
\usepackage{array}
\usepackage{enumitem}
\usepackage{titlesec}
\usepackage{hyperref}
\hypersetup{colorlinks=true, linkcolor=black, urlcolor=blue}

% headings: bold, barely larger than the text, well spaced
\titleformat{\section}{\normalfont\large\bfseries}{\thesection.}{0.6em}{}
\titleformat{\subsection}{\normalfont\normalsize\bfseries}{\thesubsection.}{0.6em}{}
\titlespacing*{\section}{0pt}{1.5em}{0.5em}
\titlespacing*{\subsection}{0pt}{1.2em}{0.4em}

\setlength{\parskip}{0.45em}
\setlength{\parindent}{0pt}
\newcommand{\unit}[1]{\,\mathrm{#1}}

\begin{document}

\begin{center}
  {\large\bfseries A surrogate model for Li/Mg selectivity in selective electrodialysis}\\[0.35em]
  {\normalsize Progress report}\\[0.9em]
  {\small Suresh Hemal \quad\textbullet\quad 28 September 2026}
\end{center}
\vspace{0.4em}
\hrule
\vspace{1.2em}

\section{The question being answered}

The task is to use the bench dataset to predict lithium transport with a polynomial surrogate
model, following the method of Section 2.5 of the GO-10 manuscript, and so to answer a design
question:

\begin{quote}
\emph{What operating setup gives the best lithium selectivity over magnesium, without
breaking the membranes and without leaving most of the lithium behind?}
\end{quote}

Stated as an optimisation, the quantity to be maximised is permselectivity,

\begin{equation}
P_{\mathrm{Li/Mg}} \;=\; \frac{t_{\mathrm{Li}}/t_{\mathrm{Mg}}}{c_{\mathrm{Li}}/c_{\mathrm{Mg}}},
\qquad
t_{\mathrm{Li}} = \frac{n_{\mathrm{Li}}}{n_{\mathrm{Li}} + 2n_{\mathrm{Mg}}},
\qquad
t_{\mathrm{Mg}} = \frac{2n_{\mathrm{Mg}}}{n_{\mathrm{Li}} + 2n_{\mathrm{Mg}}},
\end{equation}

where $n_i$ are the moles of each ion gained by the concentrate over a run and $c_i$ the molar
concentrations measured in the diluate at its start. $P = 1$ means the membrane transports the
two ions in whatever ratio the feed presents; $P > 1$ means lithium is favoured. Two
constraints are imposed on any answer:

\begin{itemize}[nosep]
  \item starting stack voltage no greater than $20\unit{V}$, above which membranes fail;
  \item lithium recovery of at least $35\%$, so that a selective but idle setup is rejected.
\end{itemize}

\section{The dataset}

The dataset currently holds twenty-one runs. All share a $400\unit{mg/L}$ lithium feed, ten
cell pairs, $6\unit{cm^2}$ of membrane area and $0.2\unit{L/min}$ flow. Three factors vary:

\begin{center}
\small
\begin{tabular}{@{}l l l@{}}
\toprule
Factor & Levels & Varied in \\
\midrule
Current density & 3, 6, 12, 24 mA/cm$^2$ & the current arm \\
Mg/Li ratio & 5, 10, 20, 40, 60 (by mass) & the ratio arm \\
Membrane & GO-10, pristine, PDA/PEI/15CE, & both \\
                     & PDA/GO10/PEI & \\
Starting concentrate & 0.5 g/L LiCl, 5 g/L NaCl, 5 g/L LiCl & the 24 mA/cm$^2$ runs \\
\bottomrule
\end{tabular}
\end{center}

Each run contributes a power log every 15 minutes, tank readings (conductivity, pH,
temperature and weight) for the concentrate, diluate and rinse, and five or more samples per
tank for ICP-OES analysis. In total the dataset holds 891 power readings, 1049 tank readings
and 345 samples. Seventeen runs currently carry laboratory results and can be modelled; the
remainder join the model as soon as their results are recorded, with no change to the program.

The design holds total charge constant at $48\unit{mA\,h/cm^2}$: every run at
$12\unit{mA/cm^2}$ lasted 4 hours, every run at $6$ lasted 8, at $3$ lasted 16 and at $24$
lasted 2. Current density and run duration are therefore not independent, a point returned to
in Section~\ref{sec:limits}.

\subsection*{Quantities and constants}

The model will be fitted to quantities computed from the raw readings by the definitions
below, rather than read from any derived cell. Setting them out here fixes what each symbol
means for the rest of the project.

\begin{center}
\small
\begin{tabular}{@{}l l l@{}}
\toprule
Constant & Value & Note \\
\midrule
Molar mass, Li & 6.94 g/mol & see footnote \\
Molar mass, Mg & 24.305 g/mol & \\
Charge number, Li$^+$ & $z = 1$ & \\
Charge number, Mg$^{2+}$ & $z = 2$ & each mole carries two moles of charge \\
Faraday constant & 96\,485 C/mol & \\
Membrane area & 6 cm$^2$ & constant across the series \\
Cell pairs & 10 & \\
Flow rate & 0.2 L/min & \\
Solution density & 1 g/mL & assumed, to turn a tank weight into a volume \\
\bottomrule
\end{tabular}
\end{center}

For a tank of tare weight $w_0$ at a reading $w$, and a measured concentration $c$ in mg/L of
an ion of molar mass $M$:

\begin{align}
V &= \frac{w - w_0}{1000} \ \ \text{[L]},
&
n &= \frac{c\,V}{1000\,M} \ \ \text{[mol]},
&
\text{equivalents} &= z\,n .
\end{align}

Writing $n_i$ for the moles of ion $i$ \emph{gained by the concentrate} between the first and
last sampled moment of a run, and $c_i^{\,\mathrm{dil}}$ for the molar concentration measured
in the diluate at the first of those moments:

\begin{align}
t_{\mathrm{Li}} &= \frac{n_{\mathrm{Li}}}{n_{\mathrm{Li}} + 2n_{\mathrm{Mg}}},
&
t_{\mathrm{Mg}} &= \frac{2n_{\mathrm{Mg}}}{n_{\mathrm{Li}} + 2n_{\mathrm{Mg}}},
&
P &= \frac{t_{\mathrm{Li}}/t_{\mathrm{Mg}}}{c_{\mathrm{Li}}^{\,\mathrm{dil}} / c_{\mathrm{Mg}}^{\,\mathrm{dil}}},
\\[0.4em]
\text{recovery} &= \frac{100\,n_{\mathrm{Li}}}{c_{\mathrm{Li}}^{\,\mathrm{dil}} V^{\,\mathrm{dil}}},
&
J_{\mathrm{Li}} &= \frac{n_{\mathrm{Li}}}{A\,\Delta t},
&
J &= \frac{\bar{I}}{A},
\\[0.4em]
Q &= \bar{I}\,\Delta t \ \ \text{[C]},
&
\eta_{\mathrm{Li}} &= \frac{100\,F\,n_{\mathrm{Li}}}{N\,Q},
&
R &= \frac{c_{\mathrm{Mg}}^{\,\mathrm{dil}}}{c_{\mathrm{Li}}^{\,\mathrm{dil}}},
\end{align}

with $A$ the membrane area, $N$ the number of cell pairs, $\bar{I}$ the mean measured current
and $\Delta t$ the elapsed time of the measurement window.

\footnotesize
\textbf{Footnote on the lithium molar mass.} 6.94 g/mol is used rather than the standard
6.941. The choice does not affect $P$ at all --- the molar mass appears in both
$t_{\mathrm{Li}}/t_{\mathrm{Mg}}$ and $c_{\mathrm{Li}}/c_{\mathrm{Mg}}$ and cancels exactly
--- and affects $t_{\mathrm{Li}}$ only in the fifth decimal and recovery by $0.01\%$. It is
kept because it is the value used at the bench.
\normalsize

\section{Work completed}

Nothing has been calculated from the data yet, and no result is reported here. The work so far
has been to put the dataset and the program on a footing where the calculations, once written,
can be trusted and repeated.

\subsection{The dataset is now machine-readable, and its provenance is checked}

The source workbook holds one sheet per run, with the readings, the derived quantities and the
results laid out together. It has been converted into four tables --- runs, power
readings, tank readings and samples --- each holding only what was set at the bench or
reported by the laboratory. No derived quantity is copied across; all of them will be computed
from the raw values.

The conversion copies values exactly as recorded and corrects nothing. To guarantee that, a
second, independent check walks the source sheets and compares every one of the
\textbf{10{,}494} values against the table built from it, reporting any disagreement. The
dataset can therefore be rebuilt at any time and shown to be a faithful copy of its source.

\subsection{Data-quality checks are automated}

Two kinds of check now run over the dataset. The first enforces structure: types, required
values, permitted ranges, and unique sample identifiers. The second examines the readings
themselves, and is built so that a mistyped or mis-transcribed value cannot reach a
calculation unnoticed. The rules are derived from each run's own numbers rather than from
fixed constants, so they apply unchanged to runs not yet performed.

\begin{itemize}[nosep]
  \item \textbf{Values that are physically impossible.} A tank cannot weigh less than the
        empty tank, nor change by more than a tenth of its contents between two readings.
        Anything breaching this is excluded from the calculation --- the affected moment is
        dropped and the run measured over the remaining window --- and reported.
  \item \textbf{Values that disagree with their own run.} A reading far from both its
        neighbours, while those neighbours agree with each other, is reported for review. This
        applies to conductivity, pH, voltage and current, which are measured often enough
        within a run to support the comparison.
  \item \textbf{Values that disagree with the second instrument.} Voltage and current are
        each recorded twice, on the supply and on a multimeter. The two normally track closely,
        so a disagreement between them is a sharp and independent test that needs no statistics.
  \item \textbf{Runs that cannot be calculated at all,} because the laboratory results they
        need have not arrived. These are named in the output
        with the reason, so that no result is ever quoted without the reader knowing how many
        runs it came from.
\end{itemize}

Every check reports; none edits the data. The dataset stays a faithful copy of its source, and
is rebuilt from it whenever the source changes.

Two devices do most of the work, and both judge a reading against its own run rather than
against any fixed expectation. Both are worth a paragraph, because the rest of the section
refers to them.

\textbf{1. A median-based score, to notice a change that is out of character.} For a column
read many times during a run, the changes between consecutive readings are collected. Two
summaries are taken from them: the typical change, and the typical distance of a change from
that typical value --- the \emph{median absolute deviation}. Both are medians, deliberately
rather than averages, because one wrong value pulls an average towards itself and so hides
inside the very spread it inflates, whereas a median does not move. A change is then measured in
multiples of that typical distance; beyond 3.5 multiples it is noticeable. Nothing about the
equipment or the tank size needs to be known, and a quiet run is judged on a finer scale than a
lively one, automatically.

Writing $x_1,\dots,x_n$ for one column's readings within a run, in time order, and
$d_i = x_{i+1} - x_i$ for the change between two consecutive readings:

\begin{align}
\tilde{d} &= \operatorname{median}\bigl(d_1,\dots,d_{n-1}\bigr),
&
\mathrm{MAD} &= \operatorname{median}\bigl(\,\lvert d_i - \tilde{d}\rvert\,\bigr),
&
z_i &= \frac{0.6745\,\bigl(d_i - \tilde{d}\bigr)}{\mathrm{MAD}},
\end{align}

and a change is noticeable when $\lvert z_i \rvert > 3.5$. The factor $0.6745$ is a conversion,
not a tuning knob: for data with the usual bell-shaped scatter the median absolute deviation
sits at about $0.6745$ of the standard deviation, so dividing by it puts $z$ on the familiar
scale where $3.5$ means ``three and a half times the ordinary scatter''. Where a column is held
deliberately constant through a run, $\mathrm{MAD}$ is zero and the score cannot be formed ---
which is why temperature and supply current are given a plain tolerance instead.

\textbf{2. A jump-and-back test, to tell a wrong reading from a real change.} Being noticeable
is not enough on its own, because a run genuinely changes --- often fastest at the end. What
marks a wrong reading is that the series jumps away and immediately jumps back, so the reading
disagrees with both its neighbours while those two agree with each other. Requiring the second
change to undo at least 70\% of the first separates the two cases. For a change $d_i$ that has
been found noticeable, the share undone by the change that follows it is

\begin{equation}
\rho_i \;=\; \frac{-\,d_{i+1}}{d_i}
\;=\;
\begin{cases}
\dfrac{\lvert d_{i+1}\rvert}{\lvert d_i \rvert} & \text{if the two changes have opposite signs},\\[0.9em]
0 & \text{otherwise (the series kept moving the same way)},
\end{cases}
\end{equation}

and the reading $x_{i+1}$ lying between them is judged wrong when $\rho_i \geq 0.7$. The two
rules applied to tank weight need no statistic at all: with $w_0$ the weight of the empty tank,
a reading is excluded when $x_i \leq w_0$, or when
$\lvert d_i \rvert > 0.10 \times \operatorname{median}(x_1,\dots,x_n)$.

The table below shows both tests on an invented series:

\begin{center}
\small
\begin{tabular}{@{}l l l@{}}
\toprule
 & A genuine change & One mistyped reading \\
\midrule
Readings & 1{,}270\quad 1{,}490\quad 1{,}750\quad 1{,}990 &
           1{,}270\quad 1{,}490\quad 14{,}900\quad 1{,}990 \\
Changes  & $+220$\quad $+260$\quad $+240$ &
           $+220$\quad $+13{,}410$\quad $-12{,}910$ \\
Noticeable? & no --- all three are ordinary & yes --- two changes stand out \\
Reversed?   & no --- every change continues upward &
              yes --- the second undoes 96\% of the first \\
Verdict     & the run is rising; nothing to report &
              the middle reading is wrong \\
\bottomrule
\end{tabular}
\end{center}

Used together, the first device finds candidates and the second discards the ones that are
simply the process moving. On this dataset that reduces the readings raised for review by
roughly an order of magnitude, which is the difference between a list someone reads and a list
someone ignores.

Table~\ref{tab:checks} sets out what is applied to each reading.

\begin{table}[htbp]
\centering
\footnotesize
\caption{Checks applied to each reading. \emph{Excluded} means the calculation may not use
that moment; the run is measured over the remaining window and the trim is reported.
\emph{Reported} means the value is raised for review and otherwise used unchanged.}
\label{tab:checks}
\begin{tabular}{@{}l p{8.8cm} l@{}}
\toprule
Reading & Check & Consequence \\
\midrule
Tank weight & At or below the weight of the empty tank & excluded \\
            & Changes by more than a tenth of the tank between two readings & excluded \\
            & Noticeable against the run's own changes, and at least 20 g & reported \\
\addlinespace
Conductivity & Noticeable against the run's own changes, and reversed by the next reading &
reported \\
\addlinespace
pH & Moves by more than a whole unit between readings & reported \\
   & Noticeable against the run's own changes, and reversed by the next reading & reported \\
\addlinespace
Temperature & Moves by more than 2 $^{\circ}$C between readings & reported \\
            & Outside 5 to 40 $^{\circ}$C & reported \\
\addlinespace
Supply voltage & Noticeable against the run's own changes, and reversed by the next reading &
reported \\
Meter voltage & As above & reported \\
Meter current & As above & reported \\
Supply current & No test of this kind: it is held constant through a run, so there is no
variation to measure against & --- \\
\addlinespace
\multicolumn{3}{@{}l}{\emph{Comparisons between the two instruments, which need no statistics:}} \\
Voltage pair & Supply and multimeter voltage differ by more than 1 V & reported \\
Current pair & Supply current, converted to mA, falls outside 0.62 to 1.6 times the
multimeter current & reported \\
\bottomrule
\end{tabular}
\end{table}

A wrong reading in a frequently measured column raises two changes --- the jump away and the
jump back --- which are one finding, belonging to the reading between them, and are reported as
one. Tank weight is the exception: it is recorded only when a sample is taken, five times in a
typical run, so a wrong value at the start or the end of a run has no neighbour on one side and
no jump back to pair with. That is why weight is judged against the tank's own size and tare
rather than against the run's variation, and why it is the only reading whose failure can
exclude anything.

\subsection{The checks are governed by stated constants, in one place}

Each rule compares against a constant chosen from the dataset rather than by taste, and all of
them live in a single file, so that revising one is a single recorded decision rather than an
edit scattered through the program.

\begin{center}
\footnotesize
\begin{tabular}{@{}l l p{8.3cm}@{}}
\toprule
Constant & Value & Reason for the value \\
\midrule
Impossible weight step & 10\% of the tank & Sampling removes tens of grams and water
transport adds tens of grams, so a change of this size cannot occur between two readings.
Expressed as a share so that it holds for every tank size. \\
Reportable weight step & 20 g & Below this a flagged step is noise: many series vary by well
under a gram between readings, which would make any change look extreme. \\
Median-based score & 3.5 & The cutoff in multiples of the typical distance described
above; the standard figure for this score. Used only on columns measured often enough within a
run to have a typical distance at all. \\
Jump-and-back share & 0.7 & How much of a jump the next change must undo for the reading
between them to be judged wrong rather than the run to have moved. \\
pH step & 1 unit & Consecutive pH readings move by hundredths. \\
Temperature step & 2 $^{\circ}$C & Temperature is held to a tenth of a degree, so its
variation within a run is nil and no statistic can be formed from it. \\
Voltage disagreement & 1 V & Supply and multimeter normally agree to a few hundredths of a
volt, so a gap of this size indicates a transcription error rather than instrument drift. \\
Current ratio & 0.62 to 1.6 & The supply reads consistently above the multimeter, so the two
are compared as a ratio rather than a difference. \\
\bottomrule
\end{tabular}
\end{center}

\subsection{The program}

The analysis is a small command-line program with an automated test suite, so that every
number in a future paper can be regenerated from the raw readings by a single command. Four
commands exist: one rebuilds the four data tables from the workbook, copying values exactly;
one reports, in plain language, anything in the data that fails a check;
one compares every value against the workbook cell it came from; and one runs the modelling.
Forty-nine automated tests cover the loading, the checks and the commands, and run on every
change.

The source is version-controlled and available at
\url{https://github.com/sureshhemal/limodel} --- currently a private repository, so access can
be granted on request. Its README documents the setup, the commands and the structure, and the
definitions in this report are also kept alongside the code so that the two cannot drift apart.

\section{The model to be fitted, and what the data can support}
\label{sec:limits}

The intended form is a second-order polynomial in the two settings that can be dialled,
fitted separately for each quantity a constraint is placed on:

\begin{equation}
\boxed{\;
\hat{y} \;=\; b_0 \;+\; b_1 J \;+\; b_2 R \;+\; b_3 J^2 \;+\; b_4 R^2 \;+\; b_5 JR
\;}
\qquad y \in \{\,P,\ \text{recovery},\ V_0\,\}
\end{equation}

where $J$ is current density in mA/cm$^2$ and $R$ the Mg/Li molar ratio of the feed. Six
coefficients are fitted per outcome, by least squares:

\begin{center}
\small
\begin{tabular}{@{}l l l@{}}
\toprule
Term & Reads as & What it would say \\
\midrule
$b_0$ & intercept & the level of the surface \\
$b_1 J$ & linear in current & whether more current helps or hurts \\
$b_2 R$ & linear in ratio & the cost of magnesium in the feed \\
$b_3 J^2$ & curvature in current & that a best current exists, rather than an extreme \\
$b_4 R^2$ & curvature in ratio & whether the magnesium penalty eases or worsens \\
$b_5 JR$ & interaction & whether the best current depends on the magnesium load \\
\bottomrule
\end{tabular}
\end{center}

Three quantities are fitted because the question is constrained: $P$ is maximised, while
recovery and starting voltage are the gates, and a gate cannot be checked without being
predicted. The membrane enters as a category rather than a continuous variable, so the surface
is fitted per membrane. The constants of the optimisation, as distinct from the fitted
coefficients, are:

\begin{center}
\small
\begin{tabular}{@{}l l l@{}}
\toprule
Constant & Value & Source \\
\midrule
Maximum starting voltage & 20 V & above this the membranes fail \\
Minimum lithium recovery & 35\% & provisional; may be revised \\
Current density searched & 3 to 24 mA/cm$^2$ & the tested range; no extrapolation \\
Mg/Li ratio searched & the tested range & set by the runs performed \\
\bottomrule
\end{tabular}
\end{center}

Four points follow from the design rather than from the method, and are better stated now than
discovered later.

\textbf{The model must be tested by holding out one run at a time, not by splitting the runs
into two piles.} The usual practice is to fit on three quarters of the data and test on the
remaining quarter. With seventeen runs that would mean fitting six coefficients to thirteen runs
and judging the result on four --- and the verdict would depend on which four. Holding back one
run at a time instead, and repeating until each has been held back once, uses every run for
fitting and still tests every run. As further runs are recorded, each is also a prediction the
model made before seeing it, which is a stronger test than any division of the runs already in
hand.

\textbf{The interaction term is not estimable from the present design.} Current density varies
only at the lowest Mg/Li ratio, and the ratio varies only at $12\unit{mA/cm^2}$. No run
combines a high current with a high magnesium load, so $b_5$ would be fitted from no evidence.
Two or three runs off these axes --- for example $6$ and $24\unit{mA/cm^2}$ at a middle ratio
--- would make it estimable. Until such runs exist, the term will be reported as unsupported
rather than fitted and quoted, and the model will be presented without it.

\textbf{Current density and run duration cannot be separated.} Because total charge was held
at $48\unit{mA\,h/cm^2}$, duration is exactly $48/J$. Any prediction at a given current
density therefore carries the matching duration with it, and this will be stated wherever the
model is reported.

\textbf{A quadratic in $R$ may be the wrong shape.} A parabola has a turning point, so fitted
to a quantity that falls steadily with magnesium load it will eventually turn upward and
predict rising selectivity at high ratios, which is not physical. Substituting $\ln R$ for $R$
keeps the same six-term structure and stays monotone. Which form to prefer will be decided by
the cross-validated error once the results are computed, and I will report both.

\section{What this arrangement already gives}

Everything listed here is true of the program as it stands, not of the finished model.

\textbf{Anyone with the source workbook can rebuild the dataset and run the checks
themselves.} The rebuild is a single command, and it needs nothing from me: no manual step, no
intermediate file that has to be passed along, no knowledge of how the program works.

\textbf{Recording a new run requires no change whatever to the program.} A sheet added to the
workbook is picked up on the next rebuild. Nothing is enumerated or hard-coded, so the twenty-second
run costs exactly as much effort as it costs to perform.

\textbf{Laboratory results are taken up as soon as they exist.} A run whose analysis has not
come back is set aside and named in the output; when its numbers are recorded, the same rebuild
includes it. No separate action is needed, and nothing has to be remembered.

\textbf{The dataset holds only what was observed.} It contains what was set at the bench and
what the laboratory reported --- and no derived quantity at all. Any figure produced from it
therefore traces back through stated formulas to a weight, a concentration or a meter reading,
rather than to a cell someone had already worked out.

\textbf{That copy is checked, not trusted.} A second, independent pass compares all
\textbf{10{,}494} values against the cells they came from, so the dataset's faithfulness to its
source is demonstrable rather than assumed.

\textbf{The checks are applied identically to every run.} There are no per-run adjustments and
nothing is special-cased. This matters because it removes an entire class of error: a rule that is
copied a row too far, a duration entered by hand for one run and not another, a reference that
shifts when a sheet is edited.

\textbf{Nothing can be dropped quietly.} A reading that cannot be true blocks its own use and
says so; anything merely unusual is printed for review and otherwise left alone. Any output states
how many runs it was built from, so a figure and the evidence behind it always travel together.

\textbf{The thresholds are stated, editable and recorded.} All of them live in one file, each
beside a sentence explaining it, and changing one requires no programming. Because that file is
under version control, every such change carries a date and an author --- so if a number moves,
when and why it moved can be answered.

\textbf{The program is tested.} Forty-nine automated tests cover the reading of the data, the
checks and the commands, and they run on every change. A test, once written, is permanent: a later
edit that breaks the behaviour fails immediately instead of quietly producing a different answer.

\textbf{It is built to outlast this series of experiments.} The tables are read through a single
component, so moving from files to a database later is a change of configuration rather than of
method. Sodium, potassium and calcium already have their places in the data structures, so another
series can be carried through the same program without redesigning it. And the definitions, the
constants and the reasoning in this report are versioned in the same repository as the program
that implements them, so the description and the implementation cannot quietly diverge.

\section{Next steps}

\textbf{1. Finish preparing the data.} Set aside the runs whose laboratory results have not
arrived, apply the exclusion rules above to the readings, and choose for each run the stretch of
time over which it will be measured --- normally the whole run, and a shorter stretch where a
reading at one end cannot be used.

\textbf{2. Compute one line of results per run.} For each run: how much lithium and magnesium
crossed into the concentrate, what share of the current each of them carried, the resulting
selectivity, the fraction of the lithium recovered, and the voltage the stack started at. These
come from the raw readings by the formulas in Section 2, in code, so that they can be recomputed
at any time.

\textbf{3. Check that arithmetic by hand.} Work one run through on paper, step by step, and
confirm the program agrees. Everything afterwards rests on this chain being right, so it is worth
proving once rather than assuming.

\textbf{4. Fit the equation --- three of them.} One equation predicts selectivity from the
current density and the feed ratio; a second predicts recovery; a third predicts voltage. The
second and third are needed because a limit cannot be respected unless it can be predicted.

Each equation is then tested in a way that does not flatter it: one run is held back, the
equation is fitted to the remaining runs, and the held-back run is predicted and compared with
what actually happened. That is repeated until every run has been held back once. The result is a
single, honest figure --- how far a prediction of a \emph{new} run is likely to be out --- which
is the number that should accompany any recommendation.

Two things are tried at this step and decided by that figure rather than by preference. First,
the feed ratio is used both directly and as its logarithm, because a plain squared term must
eventually curve upward and so would predict selectivity \emph{improving} at high magnesium
content, which cannot be right. Second, a different and much more flexible kind of model --- one
built from hundreds of small rules rather than a single smooth equation --- is fitted to the same
table. If the flexible model predicts no better, that is evidence the simple equation has
captured what the data contains, which is a useful thing to be able to state.

\textbf{5. Find the best allowed setting.} With the three equations in hand, step across the
range of current densities and feed ratios that were actually tested, discard every combination
whose predicted voltage exceeds 20 V or whose predicted recovery falls below 35\%, and keep the
combination with the highest predicted selectivity. Because the magnesium content of a real
brine is given rather than chosen, the useful answer is not a single point but a recommendation
per feed ratio: for this much magnesium, run this current. Reported separately for each membrane.

\textbf{6. Produce the figures and tables} for the manuscript, each regenerated from the raw
readings by a single command.

\end{document}
